% --------------------------------------------------------------------
% 서론 섹션입니다.
% Introduction section.
%
% 이 파일에서는 연구 배경, 연구의 필요성, 연구 목적과 연구 문제를 제시합니다.
% This file presents the research background, rationale, purpose, and questions.
% --------------------------------------------------------------------

% section은 왼쪽 사이드바와 목차에 표시되는 큰 발표 단위입니다.
% A section is a major presentation unit shown in the sidebar and table of contents.
\section{서론}

% subsection은 section 안의 세부 발표 단위입니다.
% A subsection is a smaller unit within a section.
\subsection{연구 배경}

% [t] 옵션은 슬라이드 내용을 위쪽부터 배치합니다.
% The [t] option aligns slide content from the top.
\begin{frame}[t]{연구 배경}

% 슬라이드 본문을 small 크기로 줄여 발표 화면에 안정적으로 맞춥니다.
% Use small text so the slide content fits comfortably.
\small

% block은 일반 설명이나 핵심 개념을 넣는 기본 상자입니다.
% A block is the default box for explanation or key concepts.
\begin{block}{문제 상황}

% itemize는 핵심 내용을 짧은 목록으로 제시할 때 사용합니다.
% Use itemize for concise bullet points.
\begin{itemize}

% tightitems는 목록 간격을 줄이는 사용자 정의 명령입니다.
% tightitems is a helper command that reduces list spacing.
\tightitems

% 아래 항목들은 예시 문장입니다. 자신의 연구 맥락에 맞게 바꿉니다.
% The following bullets are examples. Replace them with your own research context.
\item 연구 분야에서 해결해야 할 교육적, 과학적, 사회적 문제를 제시합니다.
\item 기존 연구나 현장 실천에서 충분히 설명되지 않은 지점을 정리합니다.
\item 연구 대상, 맥락, 핵심 개념을 한 문장으로 연결합니다.
\end{itemize}
\end{block}

% exampleblock은 작성 안내, 장점, 해결 방향처럼 긍정적 안내를 강조할 때 사용합니다.
% Use exampleblock to highlight writing guidance, strengths, or possible solutions.
\begin{exampleblock}{작성 팁}
이 슬라이드는 ``왜 이 연구가 필요한가?''에 답합니다. 배경 설명은 넓게 시작하되 마지막 문장은 자신의 연구 문제로 좁혀 줍니다.
\end{exampleblock}

% slidenote는 발표자 노트입니다. 기본 설정에서는 청중용 PDF에 보이지 않습니다.
% slidenote is a speaker note. It is hidden from the audience PDF by default.
\slidenote{
배경 설명에서는 너무 많은 선행 연구를 나열하기보다, 청중이 연구의 출발점을 빠르게 이해하도록 핵심 문제를 분명하게 말합니다.
}
\end{frame}

% 연구의 필요성에서는 기존 한계와 본 연구의 대응을 나란히 보여 줍니다.
% The rationale slide compares current limitations with the proposed response.
\subsection{연구의 필요성}
\begin{frame}[t]{연구의 필요성}
\small

% columns 환경은 한 슬라이드를 여러 세로 영역으로 나눕니다.
% The columns environment splits one slide into vertical regions.
\begin{columns}[T]

% 첫 번째 열입니다. 전체 본문 너비의 48%를 차지합니다.
% First column. It takes 48% of the text width.
\begin{column}{0.48\textwidth}
\begin{block}{현재 한계}
\begin{itemize}
\tightitems

% 논문 발표에서는 한계를 너무 추상적으로 쓰지 말고 구체적인 공백으로 씁니다.
% In a thesis or dissertation presentation, describe limitations as concrete gaps rather than vague problems.
\item 한계 1: 이론적 공백
\item 한계 2: 방법론적 공백
\item 한계 3: 교육 현장 적용의 어려움
\end{itemize}
\end{block}
\end{column}

% 두 번째 열입니다. alertblock은 문제 해결의 핵심 대응을 강조할 때 적합합니다.
% Second column. alertblock is useful for emphasizing the core response.
\begin{column}{0.48\textwidth}
\begin{alertblock}{본 연구의 대응}
\begin{itemize}
\tightitems

% 각 대응은 왼쪽 한계와 대응되도록 작성하면 설득력이 높아집니다.
% Each response is stronger when it directly answers a limitation in the left column.
\item 새로운 분석 관점 또는 설계 원리 제안
\item 실제 자료, 수업, 사례, 실험 등을 통한 검증
\item 연구 결과의 교육적 활용 가능성 제시
\end{itemize}
\end{alertblock}
\end{column}
\end{columns}

\vspace{0.25cm}

% source 명령은 작은 글씨로 출처나 인용 안내를 표시합니다.
% The source command displays a small source or citation note.
\source{한국어 문헌은 \texttt{\textbackslash kparencite\{kor\_key\}}, 영어 문헌은 \texttt{\textbackslash parencite\{eng\_key\}}로 인용합니다. 예: (\kparencite{kor_seo2009science}; \parencite{eng_nrc2012framework})}

\slidenote{
필요성은 연구자의 관심이 아니라 학문적, 교육적 필요로 설명합니다. 한계와 대응을 나란히 보여 주면 심사자가 연구의 위치를 빨리 파악합니다.
}
\end{frame}

% 연구 목적과 연구 문제는 뒤의 방법, 결과, 결론 전체를 이끄는 중심 슬라이드입니다.
% The purpose and research questions guide the later methods, results, and conclusion.
\subsection{연구 목적과 연구 문제}
\begin{frame}[t]{연구 목적과 연구 문제}
\small

% 연구 목적은 한 문장으로 연구 대상, 처치/분석, 밝힐 내용을 연결합니다.
% The purpose statement links the target, treatment/analysis, and expected explanation in one sentence.
\begin{block}{연구 목적}
본 연구의 목적은 \textbf{[핵심 대상]}을 중심으로 \textbf{[핵심 처치/분석/개발]}을 수행하고, 그 결과 \textbf{[밝히려는 변화 또는 원리]}를 설명하는 데 있다.
\end{block}

% 연구 문제는 번호 목록으로 제시하면 질의응답 때 참조하기 쉽습니다.
% Numbered research questions are easy to reference during Q&A.
\begin{exampleblock}{연구 문제}
\begin{enumerate}
\tightitems

% 연구 문제 수는 실제 연구 범위에 맞게 줄이거나 늘립니다.
% Adjust the number of research questions to match the real research scope.
\item 연구 문제 1: 어떤 특성, 요구, 구조가 확인되는가?
\item 연구 문제 2: 어떤 절차와 원리로 자료, 수업, 모형, 도구를 설계할 수 있는가?
\item 연구 문제 3: 적용 또는 분석 결과 어떤 변화, 효과, 의미가 나타나는가?
\end{enumerate}
\end{exampleblock}

\slidenote{
연구 문제는 뒤의 연구 방법과 1:1로 대응되어야 합니다. 각 연구 문제에 어떤 자료와 분석 방법으로 답할지 자연스럽게 이어지도록 설명합니다.
}
\end{frame}
